% Generated draft: requires mathematical and source review.
\documentclass[11pt]{article}
\usepackage[margin=25mm]{geometry}
\usepackage{fontspec}
\setmainfont{Latin Modern Roman}
\usepackage{amsmath,amssymb,mathtools,unicode-math}
\setmathfont{Latin Modern Math}
\usepackage{enumitem,needspace,xurl,hyperref}
\setlength{\parindent}{0pt}
\setlength{\parskip}{5pt}
\setlength{\emergencystretch}{3em}
\newcommand{\N}{\mathbb{N}}
\newcommand{\Z}{\mathbb{Z}}
\newcommand{\Q}{\mathbb{Q}}
\newcommand{\R}{\mathbb{R}}
\newcommand{\C}{\mathbb{C}}
\newcommand{\F}{\mathbb{F}}
\newcommand{\A}{\mathbb{A}}
\newcommand{\PP}{\mathbb{P}}
\newcommand{\E}{\mathbb{E}}
\newcommand{\eps}{\varepsilon}
\newcommand{\dd}{\,\mathrm{d}}
\begin{document}
\section*{UnsolvedMath ID: 30001677}
\textbf{Source problem code:} catalogue entry
\subsection{Definitions and mathematical statement}
Let \(X\) be a Polish space, meaning a separable topological space whose topology admits a complete metric. Let \(E\subseteq X\times X\) be a Borel equivalence relation: \(E\) is Borel in the product topology and is reflexive, symmetric, and transitive. Let \(I\subseteq\mathcal P(X)\) be a \(\sigma\)-ideal, meaning that \(\varnothing\in I\), every subset of a member of \(I\) belongs to \(I\), and every countable union of members of \(I\) belongs to \(I\). A set \(B\subseteq X\) is \(I\)-positive if \(B\notin I\). For such a set, define
\[
E\upharpoonright B=E\cap(B\times B).
\]
The supplied question asks whether, for an arbitrary Polish space \(X\), Borel equivalence relation \(E\) on \(X\), and \(\sigma\)-ideal \(I\) on \(X\), there exists a Borel \(I\)-positive set \(B\subseteq X\) such that \(E\upharpoonright B\) is ``significantly simpler'' than \(E\) on \(X\).

This is an unresolved formulation rather than a fully specified mathematical assertion: the source supplies neither a definition of ``significantly simpler'' nor a comparison relation or target class of equivalence relations. It also does not explicitly require \(I\) to be proper, that is, \(X\notin I\). If \(X\in I\), downward closure precludes any \(I\)-positive subset. No properness assumption, additional restriction on \(I\), or interpretation of simplification is imposed in this draft.
\subsection{Short English statement}
The question asks whether every Borel equivalence relation can be simplified by restricting it to a Borel set that is positive for a given \(\sigma\)-ideal. The intended meaning of simplification and the necessary hypotheses on the ideal are unspecified, so a precise quantified target still requires clarification.
\subsection{Sources}
\begin{itemize}
\item Local registry: \texttt{data/unsolvedmath/problems.json}; supplied record with permanent UnsolvedMath ID \(30001677\), source problem code \texttt{catalogue entry}, title ``Simplifying Borel Equivalence Relations on Positive Sets.'' The record assigns the problem to the Oberwolfach Reports Open Problems collection and gives the proposed year as 2011; no individual proposer is supplied.
\item Supplied report citation and literature source labeled ``Oberwolfach problem statement'': \emph{Set Theory (2011)}, Oberwolfach Reports, pp.\ 85--140, DOI 10.4171/owr/2011/02. \url{https://doi.org/10.4171/owr/2011/02}. The supplied background describes these reports as documenting mathematical workshops and their open-problem discussions.
\item Supplied literature source labeled ``canonization on perfect cubes'': \url{https://arxiv.org/abs/1804.05174}. No authors, full bibliographic citation, or specific theorem connecting this reference to the question are supplied.
\item Literature assessment dated exactly \texttt{2026-08-21}: the supplied curation labels the problem open and states that canonization theorems give positive results for important subclasses, but that no general theorem answers the stated simplification question for arbitrary \(I\) and \(E\). This is supplied dated curation derived from the OWR material, not an independently verified literature assessment. The record expressly requests independent verification of time-sensitive status and cited claims.
\item The supplied statement-status label is \texttt{corrected\_verified}; its accompanying verification note states only that TeX notation was enclosed in math delimiters without changing the mathematical content. This does not resolve the missing mathematical specifications.
\end{itemize}
\end{document}
